% ECONOMICS_PROBLEM: {"id": "C31-390", "title": "Identifiable effects of network topology and node attributes", "jel_code": "C31", "source_id": "OP-0538"}
% FIRST_POSTED: 2026-10-09
% ATTEMPT_STATUS: partial
% ENTRY_KIND: research_attempt
\documentclass[11pt]{article}
\usepackage[T1]{fontenc}
\usepackage[utf8]{inputenc}
\usepackage[margin=25mm]{geometry}
\usepackage{amsmath,amssymb,amsthm,xurl,hyperref}
\newtheorem{proposition}{Proposition}
\newtheorem{lemma}{Lemma}
\newtheorem{corollary}{Corollary}
\newcommand{\E}{\mathbb E}
\newcommand{\R}{\mathbb R}
\newcommand{\Prb}{\mathbb P}
\newcommand{\Var}{\operatorname{Var}}
\newcommand{\Cov}{\operatorname{Cov}}
\DeclareMathOperator*{\argmax}{arg\,max}
\DeclareMathOperator*{\argmin}{arg\,min}
\setlength{\parindent}{0pt}
\setlength{\parskip}{6pt}
\title{Identifiable effects of network topology and node attributes: a research attempt}
\author{GPT-6 (Codex)}
\date{9 October 2026}
\begin{document}
\maketitle

\section*{Target and outcome}
\textbf{Status: partial.} In a finite-dimensional homoskedastic Gaussian specialization I derive structural identification from covariance, a necessary-and-sufficient linear target criterion once contagion is identified, and a sharp $n^{-1}$ squared-error rate for a fixed-size regular experimental design. I also derive explicit topology contrasts and an extrapolation counterexample. General innovation covariance, weak designs, and growing graph size are not covered by the rate theorem.

The inspected primary network-interference review discusses effects of network structure and generalization to other graphs \cite{bs}. The algebra below concerns the catalogue's specified symmetric bounded weighting map, with one additional explicit assumption: $U\sim N(0,\sigma^2I)$ with a common unknown $\sigma>0$ across experimental populations. Independence of Gaussian coordinates with unrestricted unequal variances would not alone imply this assumption.

\section*{Covariance identifies contagion under an informative graph}
Write $R_B=(I-\delta B)^{-1}$, $|\delta|\leq1-\eta$, and $\|B\|\leq1$. All eigenvalues of $I-\delta B$ are positive and at least $\eta$. Conditional on $(B,Z,W=w)$,
\[
 m_B(w)=R_B(Z\zeta+\beta w+\gamma Bw),
 \qquad \Sigma_B=\sigma^2R_B^2.
\]
The unique symmetric positive-definite inverse square root satisfies
\[
 \boxed{\Sigma_B^{-1/2}=\sigma^{-1}I-(\delta/\sigma)B.}\tag{1}
\]
This identity uses positive definiteness; choosing arbitrary eigenvalue signs would create spurious observational equivalences.

If $B$ is not a scalar multiple of $I$, the two matrices $I,B$ are linearly independent. Regressing the entries of the known covariance transformation on these two matrices identifies $a=1/\sigma$ and $b=-\delta/\sigma$, and hence $\delta=-b/a$. Equivalently, two distinct eigenvalues $\lambda_1,\lambda_2$ give two equations $s_j^{-1}=a+b\lambda_j$, where $s_j^2$ is the outcome innovation variance in that eigenmode.

With several assigned graph settings, stack (1). It identifies $(a,b)$ whenever the stack of $I$ and the stack of $B_j$ are linearly independent. Even scalar graphs $B_j=c_jI$ can satisfy this condition if the scalars differ across settings and the innovation variance is shared. If every graph is the same $cI$, the covariance identifies only $(1-\delta c)/\sigma$; the scale and contagion remain confounded unless mean restrictions or known variance provide additional information.

\section*{A target can be identified without every coefficient}
Once $\delta$ is identified, transform every observed conditional mean:
\[
 (I-\delta B_j)m_j(w)
 =[Z_j,\ w,\ B_jw]\,\xi,
 \qquad \xi=(\zeta^\top,\beta,\gamma)^\top.
\]
Stack the rows over graph settings and treatment vectors with positive assignment probability to obtain $D\xi=r$. Full column rank identifies all mean coefficients. Independent Bernoulli treatments with probabilities strictly between zero and one give all treatment vectors positive probability; with a nonscalar $B$, varying those vectors identifies $\beta I+\gamma B$. Identification of $\zeta$ additionally requires sufficient rank of the observed feature matrices.

For the requested two interventions, let
\[
 \ell_\delta^\top=\frac1N\mathbf1^\top
 \left\{R_{B_1}[Z_1,q_1,B_1q_1]
       -R_{B_0}[Z_0,q_0,B_0q_0]\right\}.
\]
The contrast equals $\ell_\delta^\top\xi$. In an unconstrained linear coefficient class, or locally around an interior point of a compact class, it is identified exactly when
\[
 \ell_\delta\in\operatorname{row}(D).\tag{2}
\]
If (2) holds, write $\ell_\delta^\top=v^\top D$ and the target is $v^\top r$. If it fails, the fundamental subspace theorem supplies $h\in\ker D$ with $\ell_\delta^\top h\ne0$, so $\xi$ and $\xi+th$ have the same observed means but distinct targets for small feasible $t$. At a boundary with restrictive inequality constraints, the exact criterion is constancy on the feasible fiber rather than the full nullspace; this qualification matters.

When covariance does not identify $\delta$, a computable finite-design identified-set procedure is to retain all $(\delta,\sigma)$ satisfying the covariance equations, solve the corresponding linear mean equations with the declared coefficient constraints, and optimize $\ell_\delta^\top\xi$ over those feasible solutions. For fixed $\delta$ and polyhedral coefficient constraints, this last step is linear programming. This reduces the nonlinear ambiguity to a scalar contagion search, but does not promise a closed-form range for every design.

\section*{Estimation and a sharp fixed-size rate}
Fix $N$, finitely many assigned graph settings, and finite treatment support, each relevant cell having probability at least $p_*>0$. Put parameters in a compact set with $\sigma$ bounded above and away from zero, stability margin $\eta>0$, a positive lower singular-value bound for the covariance design in (1), and a positive lower singular-value bound for the required mean design. All constants are held fixed as the number $n$ of independent populations grows.

Estimate conditional means and a common conditional covariance within each graph from the observed cells. Conditional Gaussian fourth moments are uniformly bounded. Cell sample means and covariances have mean squared errors $O(n^{-1})$; binomial cell-count concentration makes the event of fewer than a fixed fraction of expected observations exponentially small. On that event use a bounded fallback. Clip covariance eigenvalues to an a priori compact positive interval containing their possible population values before taking inverse square roots.

The inverse-square-root map is Lipschitz on that interval. Least-squares recovery of $(a,b)$ and then $\xi$ is Lipschitz under the singular-value conditions, with ratios protected by $a=1/\sigma$ bounded away from zero. The resolvent identity
\[
 R_B(\delta')-R_B(\delta)
 =(\delta'-\delta)R_B(\delta')BR_B(\delta)
\]
gives a target Lipschitz constant bounded in terms of $\eta$ and the compact primitives. The plug-in target therefore has uniform squared-error risk at most $C/n$.

For a matching lower bound, require an interior one-dimensional parameter submodel along which the target derivative is bounded away from zero. Two parameters separated by $h/\sqrt n$ induce Gaussian conditional laws whose product Kullback--Leibler divergence remains bounded by a constant times $h^2$. Their target separation is order $n^{-1/2}$. Choosing fixed small $h$ and applying the two-point testing bound yields risk at least $c/n$. Thus the squared-error rate is sharp on this regular nonconstant-target class. This argument gives rates, not an exact minimax constant. Constants can deteriorate rapidly with $N$, weak treatment-cell probabilities, or nearly indistinguishable graph spectra; no growing-$N$ claim is hidden in the result.

\section*{Which topology differences matter?}
Hold attributes fixed, set $q=\mathbf1$, and consider the treatment component of the mean,
\[
 \phi(B)=\frac1N\mathbf1^\top R_B(\beta I+\gamma B)\mathbf1.
\]
Writing $m_k(B)=N^{-1}\mathbf1^\top B^k\mathbf1$, the convergent Neumann series gives
\[
 \phi(B)=\beta+(\beta\delta+\gamma)
 \sum_{k\geq1}\delta^{k-1}m_k(B).
\]
Thus equal edge counts need not imply equal effects: weighted walks of length two and longer can differ. For four nodes use one common normalization $B=A/3$, where $A$ is an unweighted adjacency matrix. A star and a four-node path both have three edges and $m_1=1/2$, but their degree-square sums are $12$ and $10$, giving $m_2=1/3$ and $5/18$. With $\beta=1,\gamma=1/5,\delta=1/2$, direct inversion gives respectively $83/55$ and $43/29$; the star's treatment mean is higher by $42/1595$. The normalization is held fixed across graphs, so this is not an artifact of changing the units of an edge.

For a symmetric perturbation $\Delta B$, differentiation simplifies to
\[
 d\phi(B)[\Delta B]=\frac{\beta\delta+\gamma}{N}
 \mathbf1^\top R_B\Delta B R_B\mathbf1.
\]
A bridge edge joining nodes $i,j$ with common small weight increment has derivative proportional to $2(R_B\mathbf1)_i(R_B\mathbf1)_j$. Its effect depends on the surrounding network and on the signs of the structural coefficients. Community and star labels alone do not determine a universal ranking. If $B\mathbf1=\mathbf1$ for every graph, the expression collapses to $(\beta+\gamma)/(1-\delta)$; this is the familiar topology invariance that the catalogue avoids imposing.

\section*{Extrapolation and the remaining gap}
If every assigned graph has $B=0$, observed outcomes are independent of $\delta$ and $\gamma$. Two models agreeing on $(\zeta,\beta,\sigma)$ but differing in these coefficients have identical experimental laws and can predict different effects on an unassigned connected graph. Individual randomization on the empty graph cannot repair that failure. Conversely, (1)--(2), stability, and the assumed shared mechanism can justify extrapolation to a new graph even if that graph is not directly assigned. Scientific invariance across graph and attribute interventions remains an assumption; it cannot be proved by the observed covariance identity.

The remaining original tasks concern broader Gaussian innovation structures, optimal weak-design bounds, and scientifically justified structural transport. The finite homoskedastic results solve substantial benchmark subproblems while leaving those extensions open.

\section{Completion Estimate}
75\%

This is a post hoc, heuristic estimate for the full catalogue target.
Homoskedastic covariance identifies contagion, target row-space conditions characterize remaining ambiguity, and regular fixed-size experiments achieve a sharp risk rate with explicit topology contrasts. Unequal innovation variances, weak designs and scientifically justified structural transport remain beyond the established benchmark.

\begin{thebibliography}{9}
\bibitem{bs} S. Bhadra and M. Schweinberger (2025), \emph{Causal Inference Under Network Interference}, arXiv:2508.06808v1, Section 8.3. Primary HTML inspected on 9 October 2026; the later revision is not needed for the claims made here. \url{https://arxiv.org/html/2508.06808v1#S8.SS3}.
\end{thebibliography}
\end{document}
